% Options for packages loaded elsewhere
% Options for packages loaded elsewhere
\PassOptionsToPackage{unicode}{hyperref}
\PassOptionsToPackage{hyphens}{url}
\PassOptionsToPackage{dvipsnames,svgnames,x11names}{xcolor}
%
\documentclass[
]{article}
\usepackage{xcolor}
\usepackage[margin=1in]{geometry}
\usepackage{amsmath,amssymb}
\setcounter{secnumdepth}{5}
\usepackage{iftex}
\ifPDFTeX
  \usepackage[T1]{fontenc}
  \usepackage[utf8]{inputenc}
  \usepackage{textcomp} % provide euro and other symbols
\else % if luatex or xetex
  \usepackage{unicode-math} % this also loads fontspec
  \defaultfontfeatures{Scale=MatchLowercase}
  \defaultfontfeatures[\rmfamily]{Ligatures=TeX,Scale=1}
\fi
\usepackage{lmodern}
\ifPDFTeX\else
  % xetex/luatex font selection
\fi
% Use upquote if available, for straight quotes in verbatim environments
\IfFileExists{upquote.sty}{\usepackage{upquote}}{}
\IfFileExists{microtype.sty}{% use microtype if available
  \usepackage[]{microtype}
  \UseMicrotypeSet[protrusion]{basicmath} % disable protrusion for tt fonts
}{}
\makeatletter
\@ifundefined{KOMAClassName}{% if non-KOMA class
  \IfFileExists{parskip.sty}{%
    \usepackage{parskip}
  }{% else
    \setlength{\parindent}{0pt}
    \setlength{\parskip}{6pt plus 2pt minus 1pt}}
}{% if KOMA class
  \KOMAoptions{parskip=half}}
\makeatother
% Make \paragraph and \subparagraph free-standing
\makeatletter
\ifx\paragraph\undefined\else
  \let\oldparagraph\paragraph
  \renewcommand{\paragraph}{
    \@ifstar
      \xxxParagraphStar
      \xxxParagraphNoStar
  }
  \newcommand{\xxxParagraphStar}[1]{\oldparagraph*{#1}\mbox{}}
  \newcommand{\xxxParagraphNoStar}[1]{\oldparagraph{#1}\mbox{}}
\fi
\ifx\subparagraph\undefined\else
  \let\oldsubparagraph\subparagraph
  \renewcommand{\subparagraph}{
    \@ifstar
      \xxxSubParagraphStar
      \xxxSubParagraphNoStar
  }
  \newcommand{\xxxSubParagraphStar}[1]{\oldsubparagraph*{#1}\mbox{}}
  \newcommand{\xxxSubParagraphNoStar}[1]{\oldsubparagraph{#1}\mbox{}}
\fi
\makeatother


\usepackage{longtable,booktabs,array}
\usepackage{calc} % for calculating minipage widths
% Correct order of tables after \paragraph or \subparagraph
\usepackage{etoolbox}
\makeatletter
\patchcmd\longtable{\par}{\if@noskipsec\mbox{}\fi\par}{}{}
\makeatother
% Allow footnotes in longtable head/foot
\IfFileExists{footnotehyper.sty}{\usepackage{footnotehyper}}{\usepackage{footnote}}
\makesavenoteenv{longtable}
\usepackage{graphicx}
\makeatletter
\newsavebox\pandoc@box
\newcommand*\pandocbounded[1]{% scales image to fit in text height/width
  \sbox\pandoc@box{#1}%
  \Gscale@div\@tempa{\textheight}{\dimexpr\ht\pandoc@box+\dp\pandoc@box\relax}%
  \Gscale@div\@tempb{\linewidth}{\wd\pandoc@box}%
  \ifdim\@tempb\p@<\@tempa\p@\let\@tempa\@tempb\fi% select the smaller of both
  \ifdim\@tempa\p@<\p@\scalebox{\@tempa}{\usebox\pandoc@box}%
  \else\usebox{\pandoc@box}%
  \fi%
}
% Set default figure placement to htbp
\def\fps@figure{htbp}
\makeatother





\setlength{\emergencystretch}{3em} % prevent overfull lines

\providecommand{\tightlist}{%
  \setlength{\itemsep}{0pt}\setlength{\parskip}{0pt}}



 
\usepackage[]{biblatex}
\addbibresource{../main/references.bib}


\usepackage{float}
\floatplacement{figure}{H}
\usepackage{booktabs}
\usepackage{colortbl}
\usepackage[table]{xcolor}
\definecolor{light-gray}{gray}{0.9}
\renewcommand{\arraystretch}{1.2}
\let\oldtabular\tabular
\let\oldendtabular\endtabular
\renewenvironment{tabular}{\small\oldtabular}{\oldendtabular}
\setlength{\tabcolsep}{4pt}
\makeatletter
\@ifpackageloaded{caption}{}{\usepackage{caption}}
\AtBeginDocument{%
\ifdefined\contentsname
  \renewcommand*\contentsname{Table of contents}
\else
  \newcommand\contentsname{Table of contents}
\fi
\ifdefined\listfigurename
  \renewcommand*\listfigurename{List of Figures}
\else
  \newcommand\listfigurename{List of Figures}
\fi
\ifdefined\listtablename
  \renewcommand*\listtablename{List of Tables}
\else
  \newcommand\listtablename{List of Tables}
\fi
\ifdefined\figurename
  \renewcommand*\figurename{Figure}
\else
  \newcommand\figurename{Figure}
\fi
\ifdefined\tablename
  \renewcommand*\tablename{Table}
\else
  \newcommand\tablename{Table}
\fi
}
\@ifpackageloaded{float}{}{\usepackage{float}}
\floatstyle{ruled}
\@ifundefined{c@chapter}{\newfloat{codelisting}{h}{lop}}{\newfloat{codelisting}{h}{lop}[chapter]}
\floatname{codelisting}{Listing}
\newcommand*\listoflistings{\listof{codelisting}{List of Listings}}
\makeatother
\makeatletter
\makeatother
\makeatletter
\@ifpackageloaded{caption}{}{\usepackage{caption}}
\@ifpackageloaded{subcaption}{}{\usepackage{subcaption}}
\makeatother
\usepackage{bookmark}
\IfFileExists{xurl.sty}{\usepackage{xurl}}{} % add URL line breaks if available
\urlstyle{same}
\hypersetup{
  pdftitle={Supplementary Legends for: Computational Tool Choice Impacts CRISPR Spacer-Protospacer Detection},
  pdfauthor={Uri Neri; Antonio Pedro Camargo; Rick Beeloo; Brian Bushnell; Simon Roux},
  colorlinks=true,
  linkcolor={blue},
  filecolor={Maroon},
  citecolor={Blue},
  urlcolor={Blue},
  pdfcreator={LaTeX via pandoc}}


\title{Supplementary Legends for: Computational Tool Choice Impacts
CRISPR Spacer-Protospacer Detection}
\author{Uri Neri \and Antonio Pedro Camargo \and Rick Beeloo \and Brian
Bushnell \and Simon Roux}
\date{}
\begin{document}
\maketitle


\section{Supplementary Materials}\label{supplementary-materials}

\subsection{Tables}\label{tables}

\subsubsection{Supplementary Table S1: Tool Configuration
Details}\label{tbl-supp-tool-config}

\texttt{supp\_table\_S1\_tool\_config.tsv}

Parameters chosen for each tool to maximize recall on short spacers.
BLASTn settings emphasize ungapped alignments and permissive
identity/coverage to capture all valid matches; bowtie variants,
minimap2, MMseqs2, mummer4, sassy, strobealign, and x\_mapper entries
list versions and command-line templates.

\subsubsection{Supplementary Table S2: IMG/VR4 Contig Dataset Statistics
and Filtering Steps}\label{tbl-supp-contig-stats}

\texttt{supp\_table\_S2\_contig\_stats.tsv}

Summary of IMG/VR4 v1.1 high-confidence viral contigs after taxonomic
filtering and length cutoffs, including the final HQ benchmark subset
(fraction\_1).

\subsubsection{Supplementary Table S3: CRISPR Spacer Dataset Composition
and Statistics}\label{tbl-supp-spacer-stats}

\texttt{supp\_table\_S3\_spacer\_stats.tsv}

Statistics for the iPHoP spacer dataset before and after complexity
filtering, including length ranges, GC content, and removal criteria.

\subsubsection{Supplementary Table S4: Detailed Statistics for All
Datasets}\label{tbl-supp-simulation-params}

\texttt{supp\_table\_S4A\_real\_dataset\_sequence\_stats.tsv}

\textbf{A. Real datasets} --- Counts and length/GC summaries for spacers
and contigs across IMG/VR4 fractions (0.0005 through 1).

\texttt{supp\_table\_S4B\_simulated\_dataset\_sequence\_stats.tsv}

\textbf{B. Simulated datasets} --- Equivalent summaries for all
synthetic datasets used in the study.

\subsubsection{Supplementary Table S5: Simulation and job-generation
configuration}\label{tbl-supp-prepare-jobs}

\texttt{supp\_table\_S5\_prepare\_jobs.tsv}

Per-dataset group configuration for simulations and tool runs, noting
search-space sizes, distance thresholds, and tool inclusion/exclusion.

\subsubsection{Supplementary Table S6: Recall Values for Each Tool at
Different Mismatch Thresholds}\label{tbl-supp-recall}

\texttt{supp\_table\_S6\_recall\_values.tsv}

Values are counted per spacer--contig alignment using exact distance for
each metric. \texttt{n\_found} tallies tool-reported positive alignments
at a given distance, \texttt{n\_total} is the ground-truth positive
count, and recall is \texttt{n\_found\ /\ n\_total}. Subsections report
IMG/VR4 and synthetic datasets separately. Note that these are
aggregated per unique \texttt{alignment\_idx}, not per unique
\texttt{region\_idx} (i.e.~without applying the boundary tolerance and
per-region aggregations). The data and results displayed in the main
text refer to the per-region aggregate. The entire results (detailing
which alignment were detected by which tools and to which
\texttt{region\_idx} they belong) can be regenerated using the projects'
analysis notebook and the raw tool outputs (available in the Zenodo
deposit, see Data availability).

\subsubsection{Supplementary Table S7: Search-space-normalized metrics
at different hamming/edit
distances}\label{tbl-supp-normalized-threshold3}

\texttt{supp\_table\_S7\_search\_space\_normalized.tsv} Normalized
non-planned match rates (per spacer per Gbp and per Gbp\(^2\)) for each
fully synthetic simulation.

\subsubsection{Supplementary Table S8: Semi-synthetic non-planned match
rates by hamming threshold}\label{tbl-supp-semisynthetic}

\texttt{supp\_table\_S8\_semi\_synthetic\_rates\_hamming\_background.tsv}

Cumulative non-planned matches in the semi-synthetic dataset across
hamming thresholds (0--5), with rates normalized by spacer and contig
sequence space. Note: For the higher thresholds (≤4, ≤5) the counts are
should be taken as lower bounds as only some of the heuristic tools
could finish within reasonable time. Furthermore, so that some of the
tools could finish within said time limit, we optimised the parameters
to a maximum value of 3 mismatches (incidentally the hard limit of
Bowtie1).

\subsubsection{Supplementary Table(s) S9: Computational Resource
Usage}\label{tbl-supp-resource-usage}

Detailed resource usage metrics for all completed tool runs on both real
(IMG/VR4 subsampled) and simulated datasets. For each tool-dataset
combination, we report total CPU time (hours), wall-clock time (hours),
peak memory usage (GB), SLURM job ID, and allocated CPU cores. All tools
were allocated 64 threads and 512 GB RAM on identical hardware (dual AMD
EPYC 7543, see Methods the Reproducibility section in the main text).
The fraction\_X datasets represent subsampled fractions of IMG/VR4 HQ
contigs, while ns\_X\_nc\_Y datasets are synthetic benchmarks. The
ns\_3826979\_nc\_421431\_real\_baseline dataset is semi-synthetic, where
real spacers are searched against simulated contigs matched to IMG/VR4
characteristics.

\paragraph{A. Real data (IMG/VR4 subsampled fractions) --- completed
jobs}\label{a.-real-data-imgvr4-subsampled-fractions-completed-jobs}

\texttt{supp\_table\_S9A\_real\_data\_resource\_usage.csv}

\paragraph{B. Simulated datasets --- completed
jobs}\label{b.-simulated-datasets-completed-jobs}

\texttt{supp\_table\_S9B\_simulated\_data\_resource\_usage.csv}

\paragraph{C. Timed-out / failed jobs}\label{c.-timed-out-failed-jobs}

\texttt{supp\_table\_S9C\_timed\_out\_failed\_jobs.csv}

\subsection{Supplementary Figures}\label{supplementary-figures}

\subsubsection{Supplementary Figure S1: IMG/VR4 tool-vs-tool comparison
matrices}\label{fig-supp-real-matrix}

\texttt{FigureS1\_supp\_imgvr4\_fraction1\_tool\_comparison\_matrix.svg}

Pairwise tool-comparison matrices for IMG/VR4 fraction\_1 at exact
hamming distances (0--3). Each cell indicates regions detected by one
tool but not another; diagonals list per-tool totals.

\subsubsection{Supplementary Figure S2: Simulated dataset recall vs
spacer occurrence frequency}\label{fig-supp-sim-recall}

\texttt{FigureS2\_supp\_imgvr4\_high\_ins\_recall\_vs\_occurrence.svg}

High-occurrence-focused recall view for the high-insertion simulated
dataset (\texttt{ns\_500\_nc\_5000\_HIGH\_INSERTION\_RATE}) at hamming
distance ≤3.

\subsubsection{Supplementary Figure S3: IMG/VR4 fraction\_1 recall vs
occurrence frequency at hamming distance ≤3}\label{fig-supp-real-recall}

\texttt{FigureS3\_supp\_imgvr4\_fraction1\_recall\_vs\_occurrence.png}

Recall versus spacer occurrence for the real IMG/VR4 fraction\_1 dataset
at hamming distance ≤3.

\subsubsection{Supplementary Figure S4: Comparison of simulated vs real
spacer characteristics}\label{fig-supp-spacer-comparison}

\texttt{FigureS4A\_supp\_filtered\_iphop\_spacers\_attribute\_distributions\_multiplot.png}
\texttt{FigureS4B\_supp\_simulated\_spacer\_distributions.svg}

Distributions comparing simulated and real spacer features (k-mer
repeatability, entropy, base frequencies, GC, LCC). Panel A shows
filtered iPHoP spacers; Panel B shows fully simulated spacers
(ns\_100000\_nc\_20000).

\subsubsection{Supplementary Figure S5: Spacer characteristic
distributions by dataset}\label{fig-supp-spacer-distributions}

\texttt{FigureS5\_supp\_real\_data\_contig\_distributions.svg}

Spacer and contig characteristic distributions for real (IMG/VR4) and
simulated datasets, covering spacer size, mismatch counts, and
occurrence rates. Outliers were removed for fraction\_1 length
visualization.

\subsubsection{Supplementary Figure S6: Distribution of matched spacers
from fraction\_1}\label{fig-supp-occurrence-distribution}

\texttt{FigureS6\_supp\_fraction\_1\_matched\_spacers\_distributions.png}

Distributions of matched spacers (hamming ≤3) across length, number of
occurrences, and per-distance alignment counts.

\subsubsection{Supplementary Figure S7: Peak memory scaling (real and
simulated datasets)}\label{fig-supp-resource-memory}

\texttt{FigureS7\_supp\_resource\_usage\_memory\_2panel.svg}

Peak memory usage (RSS) scaling with dataset size. Axes are log-scaled;
tool identity encoded by color/shape.

Panel A: IMG/VR4 subsamples;

Panel B: simulated datasets.

\subsubsection{Supplementary Figure S8: Tool recall consistency across
dataset sizes}\label{fig-supp-sim-consistency}

\texttt{FigureS8A\_supp\_imgvr4\_across\_fractions\_recall\_consistency\_hamming\_le3.svg}
\texttt{FigureS8B\_supp\_simulated\_cross\_fraction\_consistency.svg}
\texttt{FigureS8C\_supp\_simulated\_heatmap\_dataset\_consistency\_hamming\_le3.svg}

Per-tool recall consistency at hamming distance ≤3. Panel A: real
IMG/VR4 subsamples. Panel B: synthetic datasets. Panel C: heatmap of
per-tool recall across simulated runs.

\subsubsection{Supplementary Figure S9: Per-simulation non-planned match
counts (Hamming vs Edit)}\label{fig-supp-per-sim-background}

Per-simulation comparisons of non-planned match counts under hamming
versus edit distance for four fully synthetic datasets
(\texttt{ns\_50000\_nc\_5000}, \texttt{ns\_75000\_nc\_5000},
\texttt{ns\_100000\_nc\_10000}, \texttt{ns\_100000\_nc\_20000}) at
thresholds 0--5.

\texttt{FigureS9A\_supp\_ns\_50000\_nc\_5000\_background\_counts\_foldchange.svg}
\texttt{FigureS9B\_supp\_ns\_75000\_nc\_5000\_background\_counts\_foldchange.svg}
\texttt{FigureS9C\_supp\_ns\_100000\_nc\_10000\_background\_counts\_foldchange.svg}
\texttt{FigureS9D\_supp\_ns\_100000\_nc\_20000\_background\_counts\_foldchange.svg}

Grouped bars for each simulation (\texttt{ns\_50000\_nc\_5000},
\texttt{ns\_75000\_nc\_5000}, \texttt{ns\_100000\_nc\_10000},
\texttt{ns\_100000\_nc\_20000}) comparing hamming and edit distance
non-planned match counts at thresholds 0--5.

\subsection{Supplementary Notes}\label{supplementary-notes}

\subsubsection{Supplementary Note 1: Coordinate Tolerance Matching
Example}\label{note-supp-coord-tolerance}

To illustrate the necessity for coordinate tolerance matching, consider
the following case from one of the synthetic dataset
(\texttt{ns\_100000\_nc\_10000}) where three different tools detected
the same ground truth region, but reported it with slightly different
coordinates:

\begin{longtable}[]{@{}
  >{\raggedright\arraybackslash}p{(\linewidth - 18\tabcolsep) * \real{0.1000}}
  >{\raggedright\arraybackslash}p{(\linewidth - 18\tabcolsep) * \real{0.1000}}
  >{\raggedright\arraybackslash}p{(\linewidth - 18\tabcolsep) * \real{0.1000}}
  >{\raggedleft\arraybackslash}p{(\linewidth - 18\tabcolsep) * \real{0.1000}}
  >{\raggedleft\arraybackslash}p{(\linewidth - 18\tabcolsep) * \real{0.1000}}
  >{\raggedleft\arraybackslash}p{(\linewidth - 18\tabcolsep) * \real{0.1000}}
  >{\raggedleft\arraybackslash}p{(\linewidth - 18\tabcolsep) * \real{0.1000}}
  >{\raggedright\arraybackslash}p{(\linewidth - 18\tabcolsep) * \real{0.1000}}
  >{\raggedright\arraybackslash}p{(\linewidth - 18\tabcolsep) * \real{0.1000}}
  >{\raggedright\arraybackslash}p{(\linewidth - 18\tabcolsep) * \real{0.1000}}@{}}
\toprule\noalign{}
\begin{minipage}[b]{\linewidth}\raggedright
Source
\end{minipage} & \begin{minipage}[b]{\linewidth}\raggedright
Spacer ID
\end{minipage} & \begin{minipage}[b]{\linewidth}\raggedright
Contig ID
\end{minipage} & \begin{minipage}[b]{\linewidth}\raggedleft
Query Start
\end{minipage} & \begin{minipage}[b]{\linewidth}\raggedleft
Query End
\end{minipage} & \begin{minipage}[b]{\linewidth}\raggedleft
Target Start
\end{minipage} & \begin{minipage}[b]{\linewidth}\raggedleft
Target End
\end{minipage} & \begin{minipage}[b]{\linewidth}\raggedright
Strand
\end{minipage} & \begin{minipage}[b]{\linewidth}\raggedright
Mismatches/Cost (tool-reported)
\end{minipage} & \begin{minipage}[b]{\linewidth}\raggedright
Notes
\end{minipage} \\
\midrule\noalign{}
\endhead
\bottomrule\noalign{}
\endlastfoot
Ground Truth & 1c47e31b\_spacer\_10034 & 1c47e31b\_contig\_1828 & - & -
& 5001 & 5036 & reverse & 1 & \\
BLASTn & 1c47e31b\_spacer\_10034 & 1c47e31b\_contig\_1828 & 1 & 34 &
5036 & 5003 & reverse & 0 & \\
MMseqs2 & 1c47e31b\_spacer\_10034 & 1c47e31b\_contig\_1828 & 1 & 35 &
5036 & 5002 & reverse & 1 & \\
Sassy & 1c47e31b\_spacer\_10034 & 1c47e31b\_contig\_1828 & - & - & 5000
& 5036 & reverse & 1 & CIGAR: 34=1I1= \\
\end{longtable}

In this example, BLASTn reports the ungapped alignment at contig
positions 5003-5036 (0 reported mismatches), MMseqs2 reports 5002-5036
(1 reported mismatch over the terminal base), and Sassy reports
5000-5036 (1 reported insertion, suggesting the position 5000 in the
contig matches that of the last position in the spacer). The ground
truth planned occurrence was at 5001-5036. Despite these coordinate
variations, all three tools detected the same biological match.
\textbf{Critically}, all three tools report alignments that are
acceptable matches to the ground truth occurrence, under a maximum
hamming distance of 1. Via this coordinate tolerance, we avoid
``double-counting'' when aggregating results across tools.

\subsection{References and Data
Availability}\label{references-and-data-availability}

Analysis notebooks and source code is available in the project
repository: https://github.com/UriNeri/spacer\_matching\_bench. The raw
tool results, the simulated datasets, and the fraction of real datasets,
have been uploaded to the project's zenodo deposit (DOI:
10.5281/zenodo.15171878)


\printbibliography



\end{document}
